\documentclass{article}
\usepackage{icmlko}
\usepackage{graphicx}
\usepackage{booktabs}
\usepackage{amsmath}

\title{소수점 아래의 승리 \\
  {\large 귀무 200 개를 \emph{전부} 이기고도 화면은 한 글자도 안 바뀐다 ---
   유의성과 크기를 처음으로 따로 판정한 사이클}}

\begin{document}
\maketitle

\begin{abstract}
직전 노트는 「이 자로는 순서를 못 잰다」였고, 프레시 티처가 그것을 무너뜨리며
「잰다 --- 19.2 분이다」를 내놓았다. 그런데 티처의 근거였던 순열 p 는
\textbf{20 뽑기로 얻을 수 있는 최솟값}이었고, 참 귀무 안의 뽑기 변동이 효과와 같은 크기였다.
이 노트는 \textbf{200 뽑기로 다시 쟀고}, 동시에 \textbf{여섯 사이클 만에 처음으로
채택 크기를 측정 전에 박았다}.
\claimx{\textbf{유의성은 이겼다} --- 귀무 200 개가 \textbf{전부} 진짜 아래였다.
$k=0$ $\to$ $p=(1+0)/201=\mathbf{0.004975}<0.01$.}
{진짜 $0.216488$ $\cdot$ 귀무 평균 $0.202635$ $\cdot$ SD $0.007031$ $\cdot$ $z=1.9702$.
귀무 최대 $0.214130$ 조차 진짜보다 \textbf{3.396 분} 아래다. 실패한 뽑기 0 개.}
\claimx{🔴 \textbf{그런데 크기로는 졌다} --- 효과 $\mathbf{0.013853}$ 일 $=$ \textbf{19.9 분}이
사전등록이 \emph{측정 전에} 박은 채택 크기 $\mathbf{1.0}$ 일에 \textbf{72.2 배} 모자란다.}
{근거는 통계가 아니라 \textbf{출력의 단위}다. 제품이 내는 답은 「언제」이고 간격은 정수 일이다.
실측으로 확인했다 --- 기준 팔 MAE $10.3807$ 일 $\to$ 시험 팔 $10.1642$ 일,
\textbf{하루로 반올림하면 둘 다 10 일}이다. 화면에 나오는 답이 한 글자도 안 바뀐다.
그러므로 ③ 「언제 무엇이 터질지」 칸은 \textbf{안 움직인다}.}
\claimx{🔴 \textbf{그리고 이 닫힘은 뽑기를 늘려도 안 열린다} --- 사전등록이 측정 전에 계산했다.}
{순열 p 는 참 꼬리확률로 수렴하므로 뽑기를 늘리면 $p<0.01$ 도달 확률이 \emph{내려간다}
(200 뽑기 $0.0849$ $\to$ 1000 뽑기 $0.0039$). 다시 열려면 뽑기가 아니라
\textbf{자료나 자}가 바뀌어야 한다.}
\claimx{🔴 \textbf{짝지은 신뢰구간은 이 자리에서 믿으면 안 된다} ---
참 귀무끼리 짝지어도 \textbf{20 쌍 중 7 쌍(35\%)} 이 「유의」로 갈렸다.}
{티처 \#74 가 「짝지어 재면 승」이라며 낸 $+0.017933\,[+0.008984,+0.027130]$ 은
점추정까지 재현되지만, 그 \emph{구간}은 뽑기 잡음을 안 담는다.
판정의 근거는 구간이 아니라 \textbf{순열 p} 다.}
\end{abstract}

\section{두 문장을 따로 낸다}

이 노트의 형식 자체가 결과다. 직전 여섯 사이클(910 $\cdot$ 915 $\cdot$ 918 $\cdot$ 919 $\cdot$ 922 $\cdot$ 925)의
사전등록에서 \textbf{채택 크기를 적은 것은 0 건}이다 --- 이번에 직접 세었다.
918 의 MDE $0.04$ 일과 919 의 「최소 효과 $=$ 1 행/1 격자」는 \textbf{검출 해상도이지 채택 크기가 아니다};
「잴 수 있으면 채택」은 크기 규율이 \emph{없는} 것과 같다.
그래서 유의성이 몇 번 나와도 칸이 안 움직였다.

\begin{figure}[h]\centering
\includegraphics[width=\linewidth]{figs/fig1.pdf}
\caption{(a) 순열 200 뽑기 --- 귀무의 오른쪽 꼬리가 정규보다 훨씬 짧다(초과첨도 $-1.36$).
(b) 유의성과 크기는 둘이다: $p=0.004975$ 로 이기고 $72.2$ 배 모자라서 졌다.
(c) 참 귀무끼리 짝지어도 $35\%$ 가 「유의」로 갈린다.}
\end{figure}

\section{사전등록이 미리 적은 것 --- 그리고 미리 진 것}

\begin{table}[h]\centering\small
\begin{tabular}{llrl}
\toprule
관문 & 자 & 도달 가능 폭 $\div$ 문턱 & 측정 \emph{전} 신고 \\
\midrule
G4 크기 & 개선(일) $\ge 1.0$ & \textbf{0.013} & 🔴 양성 방향으로 \textbf{검정력 0} \\
\bottomrule
\end{tabular}
\caption{사전등록 \S3 의 관문 표에서 한 줄만 옮긴다. \textbf{「이 문턱은 이 효과 크기로는
원리상 못 넘는다」가 측정 전에 적혔다} --- 결과를 보고 만든 도피구가 아니다.
그럼에도 \S4 의 판정은 지시받은 규칙 그대로 냈다.}
\end{table}

사전등록 시각은 \texttt{os.utime} 으로 만들지 않았다(직전 노트가 그렇게 해서 그 대조의
검정력을 0 으로 만들었다). 사전등록 직후 \texttt{sha256}$+$\texttt{mtime} 을 별도 산출물에 굳혔고,
$1$ ulp 문제로 정정할 때 \textbf{첫 스탬프를 지우지 않고} 두 번째를 새로 찍었다 ---
두 sha 가 나란히 남는다. 두 스탬프 모두 측정 시작(10{:}23{:}52 UTC)보다 앞선다.
🔴 한계도 적는다: 스탬프도 mtime 위에 서 있고, 옳은 고침인 \textbf{「사전등록만 먼저 커밋」}은
서브세션에 커밋 권한이 없어 못 했다.

\section{자기시험 --- 티처의 수를 내 손으로 재현했다}

S1 $\cdot$ S2 $\cdot$ S8 $\cdot$ S9 는 \textbf{비트로} 일치했다.
S3 $\cdot$ S4 $\cdot$ S6 $\cdot$ S7 은 최대 $1.11\times10^{-16}$($1$ ulp) 어긋났고 \emph{이유를 찾았다} ---
직전 노트의 수는 「두 평균의 차」이고 짝지은 부트의 점추정은 「차의 평균」이다.
수학적으로 같고 부동소수 결합순서만 다르다. 그래서 재현 문턱을
「비트로 같다」에서 「상대차 $\le 10^{-14}$」로 \emph{측정 전에} 옮겼고, 그 정정 사실을 스탬프에 남겼다.

\section{408 개 격자의 정체 --- 조항 60 이 걸렸던 자리}

직전 노트는 같은 산출물에 「홀드아웃 격자 $2{,}685$」와 「홀드아웃 격자 수 $2{,}277$」을
같은 이름으로 나란히 실었다. 뒤엣것은 \textbf{간격이 하나라도 있는 격자}다.
차이 \textbf{408 개}($15.196\%$)가 왜 간격을 하나도 안 내는지 이번에 갈라 세었다.

\begin{figure}[h]\centering
\includegraphics[width=\linewidth]{figs/fig2.pdf}
\caption{(a) $408 = 162+20+226$ --- 합이 맞는다. (b) 행으로는 순 $7{,}248$ 행이 좋아졌는데
MAE 를 하루로 반올림하면 둘 다 10 일이다. (c) 자기 적발: 사전등록의 검출력 계산이
정규 근사 위에 서 있었고 그 근사가 이 자에서 깨진다.}
\end{figure}

\begin{table}[h]\centering\small
\begin{tabular}{lr}
\toprule
홀드아웃 격자(분모) & $2{,}685$ \\
\quad 간격이 하나라도 있는 격자(사건 $\ge 2$) & $2{,}277$ \\
\quad 🔴 간격이 0 인 격자 & $\mathbf{408}$ \\
\qquad ① 관측률 $<0.95$ 라 자격 미달 & $162$ \\
\qquad ② 자격은 되는데 급증 사건 0 개 & $20$ \\
\qquad ③ 급증 사건이 정확히 1 개 & $226$ \\
\midrule
$2{,}277 - 2{,}102$ $=$ 홀드아웃 간격이 정확히 1 개인 격자 & $175$ \\
\bottomrule
\end{tabular}
\caption{덤으로 「첫간격제외」가 통째로 떨어뜨리는 격자 수도 갈렸다.
$2{,}102$ 는 $2{,}277$ 에서 $175$ 를 뺀 수다.}
\end{table}

\section{🔴 자기 적발 셋}

\begin{enumerate}
\item \textbf{사전등록의 검출력 계산이 틀렸다.} 측정 전에 「$p<0.01$ 은 검정력 $8.49\%$ $\cdot$
p 는 $0.025$ 근처」로 박았는데 실측 $0.004975$ 는 \textbf{예측 구간 밖}이었다.
이유를 세어 답했다 --- 귀무의 초과첨도 $-1.36$ $\cdot$ 표준화 최대 $1.63$ SD,
즉 오른쪽 꼬리가 정규보다 짧다. \textbf{다음부터 검출력은 정규 근사가 아니라
앞선 뽑기의 경험분포로 계산한다.}
\item \textbf{배선 W4 가 안 발화했다}(200 뽑기의 서로 다른 값 $181$).
파고드니 씨앗은 멀쩡했고 \textbf{자가 격자 위에 산다} --- 눈금이 $1/52{,}585$ 일 $=1.643$ 초,
유효 칸 $1{,}311$ 에서 생일 문제 기대 충돌 $15.2$ 쌍(실측 $19$).
\textbf{검사가 틀렸지 배선이 아니다.} 그래도 발화 회계는 $4/5$ 로 신고했다
(못 심은 셋까지 분모에 넣으면 $4/8=0.500$).
\item \textbf{연습 주행 산출물을 지웠다.} 자백이 가리키는 파일이 없어졌다.
\textbf{안 지우는 쪽이 옳았다} --- 자백이 가리키는 물건은 남겨야 한다.
\end{enumerate}

\section{그래서 무엇이 남나}

「순서는 잡힌다」와 「칸은 안 움직인다」는 \textbf{모순이 아니라 두 문장}이다.
이 자료에서 간격 예측은 기후값보다 \emph{확실히} 낫고, 그 낫기가 \textbf{19.9 분}이라
제품 화면에서는 사라진다. 남은 길은 뽑기를 늘리는 것이 아니라 \textbf{자료나 자}를 바꾸는 것이다 ---
그리고 이 표의 진짜 결손은 $\tau=1.20$ 익명 급증이라 \textbf{「무엇」 열이 없다}는 것이다.
목표는 「언제 \emph{무엇이}」인데 이 자료는 「언제」밖에 못 낸다.

\end{document}
